\documentclass[11pt]{article}
\usepackage{amsmath,amssymb,amsthm}
\usepackage[margin=1in]{geometry}
\usepackage{hyperref}

\newtheorem{theorem}{Theorem}
\newtheorem{lemma}[theorem]{Lemma}
\newtheorem{proposition}[theorem]{Proposition}
\newtheorem{corollary}[theorem]{Corollary}
\theoremstyle{remark}
\newtheorem{remark}[theorem]{Remark}

\DeclareMathOperator{\GV}{GV}
\DeclareMathOperator{\Vol}{Vol}
\newcommand{\Rhat}{\widehat{R}}
\newcommand{\HH}{\mathbb{H}}
\newcommand{\CC}{\mathbb{C}}
\newcommand{\RR}{\mathbb{R}}
\newcommand{\QQ}{\mathbb{Q}}
\newcommand{\ZZ}{\mathbb{Z}}

\title{Solution to Ramanujan Challenge Problem 3.1:\\[4pt]
An integral over knot polynomial roots expressing $\pi^2$}
\author{Xiang Huang}
\date{July 2026}

\begin{document}
\maketitle

\begin{abstract}
Let $A_{7_2}(M,L)$ be the $A$-polynomial of the knot $7_2$, let
$\alpha\approx0.349269$ and $\beta\approx0.406813$ be the endpoints specified in
the problem, and let $y(x)$ be the indicated branch. We prove
\[
\int_\alpha^\beta\Bigl(\log x\,\frac{dy}{y}-\log y\,\frac{dx}{x}\Bigr)
=\frac{4\pi^2}{85}.
\]
By Khoi's variation formula the integral is a difference of Godbillon--Vey
invariants of the endpoint representations; the $\beta$ endpoint is the maximal
Fuchsian representation of $\Sigma(2,3,17)$ and is evaluated by Brooks--Goldman.
The new ingredient is a proof that the two extended Bloch classes are
\emph{torsion}. This follows from a single structural fact: the endpoint minimal
polynomials are palindromic, so every non-real embedding of the endpoint fields
lies on the unit circle, and on the unit circle two of the four tetrahedron
shapes are real while the other two are exchanged by
$z\mapsto1/(1-\bar z)$, under which the Bloch--Wigner function changes sign.
Torsion gives rationality; the Merkurjev--Suslin torsion order supplies the
denominator bound $4080$; a $301$-digit evaluation then determines the value
uniquely.
\end{abstract}

\section{Setup}

Write $A=A_{7_2}(M,L)$ for the $A$-polynomial of $7_2$, of degree $22$ in $M$
and $5$ in $L$, as displayed in the problem statement. Put
\[
\alpha:\ A(\alpha,\alpha^{1/2})=0,\qquad
\beta:\ A(\beta,\beta)=0,
\]
the roots nearest $0.349269$ and $0.406813$, and let $y(x)$ be the branch with
$y>0$, $y'\le-2$ on $[\alpha,\beta]$. Set
\[
I:=\int_\alpha^\beta\Bigl(\log x\,\frac{dy}{y}-\log y\,\frac{dx}{x}\Bigr).
\]

It is convenient to use the eigenvalue coordinate. Writing $M=a^2$, $L=a$ at the
$\alpha$ endpoint and $M=L=b$ at the $\beta$ endpoint, and factoring
$A(a^2,a)$ resp.\ $A(x,x)$ over $\QQ$, the endpoints are algebraic numbers with
explicit minimal polynomials:

\begin{proposition}[Endpoint fields]\label{prop:fields}
The endpoint eigenvalues satisfy
\begin{align}
f_\alpha(a)&=a^{12}-3a^{11}+4a^{10}-5a^{9}+6a^{8}-7a^{7}+7a^{6}
 -7a^{5}+6a^{4}-5a^{3}+4a^{2}-3a+1,\label{eq:falpha}\\
f_\beta(b)&=b^{16}-7b^{15}+22b^{14}-48b^{13}+87b^{12}-133b^{11}+178b^{10}
 -211b^{9}+223b^{8}\notag\\
 &\qquad-211b^{7}+178b^{6}-133b^{5}+87b^{4}-48b^{3}+22b^{2}-7b+1.
 \label{eq:fbeta}
\end{align}
Both are irreducible over $\QQ$ and \emph{palindromic}. Their signatures are
$(r_1,r_2)=(2,5)$ and $(4,6)$, and
\[
a=0.59098942867025644\ldots,\qquad a^2=\alpha,\qquad
b=0.40681308133679\ldots=\beta .
\]
\end{proposition}

\section{Reduction to a regulator difference}

Khoi's variation formula for the one-form $\log x\,d\log y-\log y\,d\log x$ on
an $A$-polynomial curve identifies $I$ with a difference of Godbillon--Vey
invariants of the two endpoint representations of $\pi_1$ of the $7_2$
complement, with the normalisation
\begin{equation}\label{eq:khoi}
I=\frac{\GV(\rho_\alpha)-\GV(\rho_\beta)}{4}.
\end{equation}

\begin{proposition}[$\beta$ endpoint]\label{prop:beta}
$S^3_{-1}(7_2)\cong\Sigma(2,3,17)$, the $\beta$ representation factors through
$\pi_1(\Sigma(2,3,17))$ as the maximal Fuchsian representation, and
\[
\GV(\rho_\beta)=\frac{242\pi^2}{51}.
\]
\end{proposition}
This is Brooks--Goldman applied to the Seifert data of $\Sigma(2,3,17)$: with
$\chi_{\mathrm{orb}}=-11/102$ and $e=-1/102$, the maximal value is
$4\pi^2\chi_{\mathrm{orb}}^2/e=-242\pi^2/51$ up to orientation.

By \eqref{eq:khoi} and Proposition~\ref{prop:beta}, the conjecture
$I=4\pi^2/85$ is \emph{equivalent} to
\begin{equation}\label{eq:target}
\GV(\rho_\alpha)=\frac{242\pi^2}{51}+\frac{16\pi^2}{85}
=\frac{1258\pi^2}{255}=\frac{74\pi^2}{15}.
\end{equation}

Finally, the Neumann--Zagier differential identity
$d\mathcal R=-(\log M\,d\log L-\log L\,d\log M)$ gives
\begin{equation}\label{eq:NZ}
I=-\Delta\mathcal R,\qquad
\mathcal R=\sum_{j\in\{T,U,V,W\}}\Rhat(z_j),\qquad
\Rhat(z)=\mathrm{Li}_2(z)+\tfrac12\log z\log(1-z)-\tfrac{\pi^2}{6},
\end{equation}
where $z_T,\dots,z_W$ are the four tetrahedron shapes of the $7_2$ ideal
triangulation (SnapPy's triangulation of $7_2$ has exactly four tetrahedra, all
positively oriented at the complete structure).

\section{The deformation chart}

With $X=M^2$ put
\begin{equation}\label{eq:chart}
u=\frac{L+X^3}{X(L+X)},\qquad
r=\frac{-(1+\sqrt{1+4u^2})}{2u},\qquad
T=1-r^2,\quad U=u,\quad V=\frac uX,\quad W=\frac1{1-uX}.
\end{equation}
At the two endpoints this gives, to the displayed precision,
\[
\begin{array}{llll}
T_\alpha=-0.15787528\ldots,&U_\alpha=6.8157987\ldots,&
V_\alpha=55.872474\ldots,&W_\alpha=5.9329217\ldots,\\[2pt]
T_\beta=-0.25828783\ldots,&U_\beta=4.3429623\ldots,&
V_\beta=26.241958\ldots,&W_\beta=3.5555141\ldots,
\end{array}
\]
all real, with $T<0$ and $U,V,W>1$: the endpoint representations are real, and
the tetrahedra are flat.

The chart \eqref{eq:chart} involves $\sqrt{1+4u^2}$, so a priori the shapes lie
in a quadratic extension of the endpoint field. They do not.

\begin{lemma}[Shape field]\label{lem:shapefield}
$1+4u^2$ is a square in the endpoint field itself. Explicitly, with
$s_\alpha,s_\beta\in\ZZ[z]$
\[
\begin{aligned}
s_\alpha&=12z^{11}-30z^{10}+32z^{9}-42z^{8}+50z^{7}-58z^{6}+54z^{5}-56z^{4}
+42z^{3}-38z^{2}+28z-21,\\
s_\beta&=60z^{15}-394z^{14}+1150z^{13}-2386z^{12}+4196z^{11}-6176z^{10}
+8022z^{9}-9202z^{8}\\
&\qquad+9406z^{7}-8592z^{6}+6958z^{5}-4960z^{4}+3064z^{3}-1546z^{2}+646z-137,
\end{aligned}
\]
one has $s_\alpha(a)^2=1+4u^2$ in $\QQ(a)$ and $s_\beta(b)^2=1+4u^2$ in
$\QQ(b)$. Hence all four shapes lie in the endpoint field, and the shape field
$E$ coincides with $F$.
\end{lemma}

\begin{proof}
Writing $u$ itself as a polynomial modulo the minimal polynomial,
\[
\begin{aligned}
u_\alpha&=-6z^{11}+14z^{10}-15z^{9}+21z^{8}-23z^{7}+28z^{6}-25z^{5}+27z^{4}
-20z^{3}+18z^{2}-13z+10,\\
u_\beta&=-6z^{15}+41z^{14}-125z^{13}+266z^{12}-474z^{11}+711z^{10}-935z^{9}
+1088z^{8}\\
&\qquad-1127z^{7}+1043z^{6}-857z^{5}+620z^{4}-389z^{3}+202z^{2}-85z+21,
\end{aligned}
\]
the identities $s^2-(1+4u^2)=c\cdot f$ hold in $\ZZ[z]$ with explicit cofactors
$c$; for $\alpha$,
\[
c_\alpha=-48z^{9}+20z^{8}+12z^{7}+104z^{6}+12z^{5}+104z^{4}+12z^{3}+56z^{2}
-16z+40 .
\]
Each identity is a single polynomial multiplication. All four are verified in
Lean (\S\ref{sec:lean}).
\end{proof}

This point matters: the torsion argument below controls the embeddings of the
field over which the class is defined. Because $s_\alpha,s_\beta$ lie in $F$
itself, the conclusion transports under \emph{every} embedding simultaneously,
with no place-by-place analysis.

\section{The two classes are torsion}

This section contains the new content. Let $D$ denote the Bloch--Wigner
dilogarithm, so that the Borel regulator of the extended Bloch element
$\xi=\sum_j[\widehat{z_j}]$ at an embedding $\sigma$ is
$D_\sigma(\xi)=\sum_j D(\sigma(z_j))$. Recall
\begin{equation}\label{eq:Dsym}
D(\bar z)=-D(z),\qquad D\Bigl(\frac1{1-z}\Bigr)=D(z),\qquad
D(z)=0\ \text{for } z\in\RR .
\end{equation}

\begin{lemma}[Unit circle]\label{lem:circle}
Every non-real embedding of $\QQ(a)$ and of $\QQ(b)$ sends the endpoint
eigenvalue to a point of modulus~$1$.
\end{lemma}

\begin{proof}
$f_\alpha$ and $f_\beta$ are palindromic, so each factors as
$f(z)=z^{n/2}g(z+z^{-1})$ for a polynomial $g$ of degree $n/2$ over $\QQ$.
For \eqref{eq:falpha},
\[
g_\alpha(w)=w^6-3w^5-2w^4+10w^3-w^2-7w+1,
\]
and for \eqref{eq:fbeta},
\[
g_\beta(w)=w^8-7w^7+14w^6+w^5-25w^4+9w^3+12w^2-3w-1 .
\]
Both are \emph{totally real}: $g_\alpha$ has $6$ real roots and $g_\beta$ has
$8$. A root $z$ of $f$ lies on $|z|=1$ if and only if $z+z^{-1}\in[-2,2]$.
Now $g_\alpha$ has exactly $5$ roots in $[-2,2]$ and one outside, giving
$10$ roots of $f_\alpha$ on the unit circle and $2$ real ones; $g_\beta$ has
exactly $6$ in $[-2,2]$ and $2$ outside, giving $12$ on the circle and $4$ real.
These counts exhaust the degrees $12$ and $16$.
\end{proof}

\begin{lemma}[$u$ is real]\label{lem:ureal}
In either chart, $u(1/a)=u(a)$ identically. Consequently, at an embedding with
$|a|=1$ one has $\bar u=u$, i.e.\ $u\in\RR$.
\end{lemma}

\begin{proof}
For the $\alpha$ chart ($L=a$, $X=a^4$), \eqref{eq:chart} gives
\[
u(a)=\frac{a+a^{12}}{a^{4}(a+a^{4})}
=\frac{1+a^{11}}{a^{4}(1+a^{3})},
\]
and substituting $a\mapsto1/a$ returns the same rational function; the identity
$u(1/a)-u(a)=0$ holds in $\QQ(a)$. For the $\beta$ chart ($L=b$, $X=b^2$) one
gets $u(b)=b^{-2}(b^4-b^3+b^2-b+1)$, again invariant under $b\mapsto1/b$.
On $|a|=1$ we have $\bar a=1/a$, hence $\bar u=u(\bar a)=u(1/a)=u(a)=u$.
\end{proof}

\begin{lemma}[$T$ and $U$ are real]\label{lem:TU}
At an embedding with $|a|=1$, $T$ and $U$ are real; hence $D(T)=D(U)=0$.
\end{lemma}

\begin{proof}
$U=u$ is real by Lemma~\ref{lem:ureal}. Then $1+4u^2>0$, so $\sqrt{1+4u^2}$ is
real, $r$ is real, and $T=1-r^2$ is real. Apply \eqref{eq:Dsym}.
\end{proof}

\begin{lemma}[$V$ and $W$ cancel]\label{lem:VW}
At an embedding with $|a|=1$, $W=1/(1-\bar V)$, and hence $D(V)+D(W)=0$.
\end{lemma}

\begin{proof}
$|a|=1$ gives $|X|=1$, so $\bar X=1/X$. With $u$ real,
\[
\bar V=\overline{\Bigl(\frac uX\Bigr)}=\frac u{\bar X}=uX,
\qquad\text{hence}\qquad
W=\frac1{1-uX}=\frac1{1-\bar V}.
\]
By \eqref{eq:Dsym}, $D(W)=D\bigl(1/(1-\bar V)\bigr)=D(\bar V)=-D(V)$.
\end{proof}

\begin{lemma}[Flattening ambiguity]\label{lem:fill}
Choose flattenings $[z_j;p_j,q_j]$ satisfying Neumann's edge and parity
conditions together with the $(-1)$-filling peripheral condition, so that
$\xi_\alpha,\xi_\beta$ represent the filled representations. At both endpoints
all four shapes are \emph{real}, and then \emph{any} change of the flattening
integers alters $\mathrm{Re}[\Delta\mathcal R]/\pi^2$ by a
\textbf{half-integer}. Consequently no flattening has to be exhibited: whatever
the choice, the denominator of $\mathrm{Re}[\Delta\mathcal R]/\pi^2$ divides
$2Q=4080$.
\end{lemma}

\begin{proof}
The $(p,q)$-dependent part of Neumann's
$\widehat{\mathcal R}([z;p,q])
=\operatorname{Li}_2(z)+\tfrac12\log z\log(1-z)+\tfrac{\pi i}2\bigl(q\log z+p\log(1-z)\bigr)
-\tfrac{\pi^2}6$
is $\tfrac{\pi i}2(q\log z+p\log(1-z))$. Writing the principal branches as
$\log z=A+i\theta_1$ and $\log(1-z)=B+i\theta_2$ with $A,B$ real and
$\theta_k\in(-\pi,\pi]$, that part equals
$\tfrac{\pi i}2(qA+pB)-\tfrac\pi2(q\theta_1+p\theta_2)$, so its real part is
$-\tfrac\pi2(q\theta_1+p\theta_2)$.

At both endpoints every shape is real, so each $\theta_k$ is $0$ or $\pi$:
for $z<0$ (our $T$) one has $\theta_1=\pi,\ \theta_2=0$, and for $z>1$ (our
$U,V,W$) one has $\theta_1=0,\ \theta_2=\pi$. Hence each tetrahedron contributes
an integer multiple of $\pi^2/2$, and so does the sum over the four tetrahedra
and the difference between the endpoints. Verified numerically to $60$ digits
over a range of flattening shifts in \texttt{scripts/flattening\_ambiguity.py}
(worst deviation from a half-integer: $3\times10^{-61}$).
\end{proof}

\begin{remark}
This replaces a weaker earlier argument. Neumann's Dehn-filling core correction
is \emph{linear} in the core complex length $\lambda_{\mathrm{core}}$, not
quadratic --- the contribution to the regulator is
$\tfrac12\lambda_{\mathrm{core}}\pi i \pmod{\pi^2\ZZ}$ --- so for
$\lambda_{\mathrm{core}}=i\theta$ purely imaginary it is $-\tfrac12\theta\pi$,
which is real and \emph{not} in general a rational multiple of $\pi^2$; integral
surgery does not force $\theta/\pi\in\QQ$. The reasoning
``$\lambda_{\mathrm{core}}^2$ is real, hence no volume contribution'' is
therefore unavailable, and so is the stronger claim that the correction is
``absorbed'' into the flattening --- that would require constructing the filled
flattening explicitly. Lemma~\ref{lem:fill} avoids both: it needs only that the
ambiguity is bounded, which the reality of the shapes supplies. The $301$-digit
rationality of \eqref{eq:pin} is a further consistency check: had an irrational
$-\tfrac\pi2(\theta_\alpha-\theta_\beta)$ been silently omitted, the computed
value could not have reconstructed to a rational at all.
\end{remark}

\begin{lemma}[Orientation signs]\label{lem:signs}
For the $7_2$ triangulation used here, $\varepsilon_j=+1$ for all four
tetrahedra. In particular $\varepsilon_V=\varepsilon_W$, and the regulator is the
\emph{unsigned} sum $\sum_j D(\sigma(z_j))$.
\end{lemma}

\begin{proof}
The sign $\varepsilon_j$ is the orientation of tetrahedron $j$ relative to the
orientation of the manifold. It is combinatorial data fixed by the oriented
triangulation, and it does \emph{not} change along the deformation arc --- in
particular it must not be read off the sign of $\mathrm{Im}\,z_j$, which
degenerates exactly where the tetrahedra become flat, which is where our
endpoints sit.

SnapPy produces oriented ideal triangulations for orientable manifolds, so
$\varepsilon_j=+1$ uniformly. This is confirmed by the volume identity
$\Vol(M)=\sum_j\varepsilon_j D(z_j)$ at the complete structure: the
\emph{all-plus} sum reproduces the volume,
\[
\sum_{j=1}^{4}D(z_j)=3.331744231641\ldots,\qquad
\Vol(7_2)=3.331744231600\ldots,
\]
agreeing to $4\times10^{-11}$, and likewise for the $(-1/2)$-filling,
$\sum_j D(z_j)=2.582170632117\ldots$ against
$\Vol(S^3_{-1/2}(7_2))=2.582170632120\ldots$, agreeing to $3\times10^{-12}$. A
single sign flip would change the all-plus sum by $2D(z_j)\ge0.79$, so no
$\varepsilon_j$ can be $-1$. See \texttt{scripts/orient\_signs.py}.
\end{proof}

\begin{remark}
This step is worth isolating because it is the one place where the shape identity
$W=1/(1-\bar V)$ is \emph{not} enough. That identity gives $D(V)+D(W)=0$ for the
unsigned values; in the signed sum one gets
$(\varepsilon_V-\varepsilon_W)D(V)$, so the cancellation genuinely requires
$\varepsilon_V=\varepsilon_W$ --- a combinatorial fact about the triangulation,
not a dilogarithm identity.
\end{remark}

\begin{theorem}[Torsion]\label{thm:torsion}
The extended Bloch classes $\xi_\alpha$ and $\xi_\beta$ of the two endpoint
representations are torsion. Consequently
\[
\mathrm{Re}\,[\Delta\mathcal R]\in\pi^2\QQ .
\]
\end{theorem}

\begin{proof}
Let $\sigma$ be an embedding of the endpoint field. If $\sigma$ is real, all four
shapes are real and $D_\sigma=0$ by \eqref{eq:Dsym}. If $\sigma$ is non-real,
then $|\sigma(a)|=1$ by Lemma~\ref{lem:circle}, so by
Lemmas~\ref{lem:TU}--\ref{lem:signs},
\[
D_\sigma(\xi)=D(T)+D(U)+D(V)+D(W)=0+0+D(V)-D(V)=0,
\]
and by Lemma~\ref{lem:fill} the passage from the relative class to the closed
class adds nothing to the Bloch--Wigner value. Thus the Borel regulator of
$\xi$ vanishes at every complex embedding, and a Bloch element whose Borel
regulator vanishes at all complex embeddings is torsion. The Rogers regulator of
a torsion class lies in $\pi^2\QQ$.
\end{proof}

\begin{remark}
Numerically, $|D_\sigma|<10^{-60}$ at each of the five conjugate pairs of
$\QQ(a)$; Theorem~\ref{thm:torsion} explains this as an identity rather than a
coincidence. We also checked that no embedding of $\QQ(a)$ gives
$D_\sigma=\pm\Vol(S^3_{-1/2}(7_2))=\pm2.58217063\ldots$, i.e.\ the $\alpha$
component does not carry the discrete faithful character. (The image of
$\rho_\alpha$ is in fact \emph{non-discrete}: the trace of a commutator word is
an algebraic integer of degree $6$ that is not of the form
$\zeta+\zeta^{-1}$ for any root of unity, so an elliptic element has infinite
order. Non-discreteness does not obstruct torsion of the Bloch class.)
\end{remark}

\section{Denominator bound and conclusion}

\begin{proposition}[Denominator]\label{prop:denom}
Let $F_\alpha=\QQ(a)$ and $F_\beta=\QQ(b)$. Then
\[
w_2(F_\alpha)=120=2\cdot2^2\cdot3\cdot5,\qquad
w_2(F_\beta)=408=2\cdot2^2\cdot3\cdot17,
\]
with $\operatorname{disc}F_\alpha=-5^6\cdot59\cdot3881^2$ and
$\operatorname{disc}F_\beta=17^{14}\cdot1597$. By the identification of the
extended Bloch group with $K_3^{\mathrm{ind}}$ and the Merkurjev--Suslin torsion
formula $|K_3^{\mathrm{ind}}(F)_{\mathrm{tors}}|=w_2(F)$, the order of each
class divides the corresponding $w_2$; together with the factor $2$ from
Lemma~\ref{lem:fill} the denominator of
$\mathrm{Re}[\Delta\mathcal R]/\pi^2$ divides
\[
Q_0=\operatorname{lcm}(120,408)=2040,\qquad Q=2Q_0=4080,
\]
\end{proposition}

\begin{theorem}[Main]\label{thm:main}
$\displaystyle
\int_\alpha^\beta\Bigl(\log x\,\frac{dy}{y}-\log y\,\frac{dx}{x}\Bigr)
=\frac{4\pi^2}{85}.$
\end{theorem}

\begin{proof}
By Theorem~\ref{thm:torsion} and Proposition~\ref{prop:denom},
$\mathrm{Re}[\Delta\mathcal R]/\pi^2$ is a rational number with denominator
dividing $Q=4080$. Two distinct rationals with denominator at most $Q$ differ by
at least $1/Q^2$, so such a rational is determined uniquely by any approximation
of error less than $1/(2Q^2)=3.00\times10^{-8}$.

Evaluating \eqref{eq:NZ} at $1000$ bits of working precision from the exact
endpoint data gives
\[
\frac{\mathrm{Re}[\Delta\mathcal R]}{\pi^2}
=-0.0470588235294117647058823529411764705882352941176470588235294\ldots
\]
with
$\bigl|\mathrm{Re}[\Delta\mathcal R]/\pi^2+\tfrac4{85}\bigr|
=1.63\times10^{-301}$,
far below the separation threshold. Rational reconstruction with denominator
bound $Q$ therefore returns \emph{uniquely}
\begin{equation}\label{eq:pin}
\frac{\mathrm{Re}[\Delta\mathcal R]}{\pi^2}=-\frac4{85}.
\end{equation}
(Consistency: $85\mid4080$, with quotient $48$.) By \eqref{eq:NZ},
$I=-\Delta\mathcal R=4\pi^2/85$.
\end{proof}

\begin{corollary}
$\GV(\rho_\alpha)=\dfrac{74\pi^2}{15}$, by \eqref{eq:target}.
\end{corollary}

\begin{remark}[Robustness of the denominator bound]\label{rem:robust}
The passage from ``$\xi$ is torsion of order dividing $m$'' to ``the Rogers value
lies in $\tfrac1m\pi^2\ZZ$'' is normalisation-sensitive: with Neumann's
$\widehat{\mathcal R}$, whose period lattice is $\pi^2\ZZ$, the denominator
divides $m=w_2(F)$ exactly as used above, but other normalisations of the
Cheeger--Chern--Simons class carry conventional factors of $6$, $12$ or $24$.

None of this affects the conclusion. The numerical certificate has error
$1.63\times10^{-301}$ while the separation of rationals of denominator at most $Q$
is $1/Q^2$, so \eqref{eq:pin} follows from \emph{any} denominator bound
$85\le Q\le 10^{150}$ (the sharp ceiling is $\sqrt{1/(2\varepsilon)}=1.75\times10^{150}$) --- roughly $146$ orders of magnitude of slack. A stray
factor of $6$ or $24$, or indeed of $10^{146}$, changes nothing. The Lean theorem
\texttt{regulator\_quotient\_eq\_robust} states exactly this, and is instantiated
there at the inflated bound $Q=24\cdot2040$, comfortably above the $4080$ used here.
\end{remark}

\section{Lean verification}\label{sec:lean}

The directory \texttt{lean/} is a standalone Lake project pinned to Lean
\texttt{v4.29.0} and Mathlib \texttt{v4.29.0}. It builds with
\texttt{lake exe cache get} followed by \texttt{lake} \texttt{build}
(\texttt{Build completed successfully (3295 jobs)}), contains no \texttt{sorry}
and no \texttt{native\_decide}, and every one of the $37$ audited declarations
depends only on \texttt{[propext, Classical.choice, Quot.sound]}.

\begin{center}
\begin{tabular}{ll}
\hline
Lemma~\ref{lem:circle} & \texttt{TraceRoots.gAlpha\_totally\_real},
\texttt{gBeta\_totally\_real};\\
& \texttt{UnitCircle.norm\_eq\_one\_of\_trace\_real\_abs\_le\_two}\\
Lemma~\ref{lem:ureal} & \texttt{ChartSymmetry.chartUAlpha\_isReal\_of\_norm\_one}\\
Lemma~\ref{lem:shapefield} & \texttt{ShapeField.sPolyAlpha\_sq},
\texttt{ShapeField.sPolyBeta\_sq}\\
Lemmas~\ref{lem:TU}--\ref{lem:VW} &
\texttt{ShapeCancellation.four\_shape\_sum\_eq\_zero}\\
Lemmas~\ref{lem:circle}--\ref{lem:VW} combined &
\texttt{MainTheorem.four\_shape\_sum\_vanishes\_of\_trace\_real}\\
Theorem~\ref{thm:main}, final step & \texttt{MainTheorem.regulator\_value}\\
Remark~\ref{rem:robust} & \texttt{regulator\_quotient\_eq\_robust}\\
\hline
\end{tabular}
\end{center}

Three inputs are stated as hypotheses rather than proved in Lean, and are cited
above: the three Bloch--Wigner functional equations \eqref{eq:Dsym} (bundled as
the structure \texttt{BlochWignerLaws}); the Merkurjev--Suslin denominator bound
of Proposition~\ref{prop:denom}; and the $301$-digit numerical evaluation. The
first is standard; the second enters only through Remark~\ref{rem:robust}, which
needs no more than its existence; the third is reproduced by \texttt{scripts/}.

\section{Reproducing the computations}

All computations are in \texttt{scripts/}; each is standalone
(Sage, or Python with \texttt{mpmath}/\texttt{SnapPy}).

\begin{center}
\begin{tabular}{ll}
\hline
\texttt{falpha2.sage} & derives \eqref{eq:falpha} from $A(a^2,a)$\\
\texttt{confirm.sage} & Lemma~\ref{lem:circle}: unit-circle counts via $g_\alpha,g_\beta$\\
\texttt{exact\_proof.sage} & Lemma~\ref{lem:ureal}: the identity $u(1/a)=u(a)$\\
\texttt{mechanism.py} & Lemmas~\ref{lem:TU}--\ref{lem:VW} at every embedding\\
\texttt{flattening\_ambiguity.py} & Lemma~\ref{lem:fill}: the ambiguity is a half-integer\\
\texttt{shapefield.sage} & Lemma~\ref{lem:shapefield}: the square-root certificates\\
\texttt{orient\_signs.py} & Lemma~\ref{lem:signs}: all-plus sum reproduces the volume\\
\texttt{alpha\_component.py} & no embedding realises $\Vol(S^3_{-1/2}(7_2))$\\
\texttt{disc\_test.py} & non-discreteness (cyclotomic trace test)\\
\texttt{w2bound.sage} & Proposition~\ref{prop:denom}: $w_2=120,408$\\
\texttt{shapes300.sage}, \texttt{pin.sage} & the $301$-digit evaluation and reconstruction\\
\hline
\end{tabular}
\end{center}

\begin{thebibliography}{9}
\bibitem{Khoi} V.~T.~Khoi, \emph{On the integral of $\log x\,dy/y-\log y\,dx/x$
over the $A$-polynomial curves}, Acta Math.\ Vietnam.\ \textbf{33} (2008),
519--528; arXiv:0811.2725. Theorem~2.1.
\bibitem{BG} R.~Brooks and W.~Goldman, \emph{Volumes in Seifert space},
Duke Math.\ J.\ \textbf{51} (1984), 529--545.
\bibitem{Neumann} W.~D.~Neumann, \emph{Extended Bloch group and the
Cheeger--Chern--Simons class}, Geom.\ Topol.\ \textbf{8} (2004), 413--474;
Theorems 12.1 and 14.5.
\bibitem{Zickert} C.~K.~Zickert, \emph{The extended Bloch group and algebraic
$K$-theory}, J.\ reine angew.\ Math.\ \textbf{704} (2015), 21--54; Theorem~1.1.
\bibitem{Franc} S.~Francaviglia, \emph{Hyperbolic volume of representations of
fundamental groups of cusped 3-manifolds}, IMRN 2004, no.~9, 425--459.
\bibitem{BW} M.~Brittenham and Y.-Q.~Wu, \emph{The classification of exceptional
Dehn surgeries on 2-bridge knots}, Comm.\ Anal.\ Geom.\ \textbf{9} (2001),
97--113.
\bibitem{CCGLS} D.~Cooper, M.~Culler, H.~Gillet, D.~D.~Long, P.~B.~Shalen,
\emph{Plane curves associated to character varieties of 3-manifolds},
Invent.\ Math.\ \textbf{118} (1994), 47--84.
\end{thebibliography}

\end{document}
